% Source: https://artofproblemsolving.com/community/c3h1709569_junior_regional_olympiad__fbh_2018_grade_7_problem_2
% Stable identity: aops-topic-1709569-q01
Let the number be $\overline{abc}$ with digits $a,b,c\in\{1,\dots,9\}$. Crossing out one digit gives $\overline{bc}$, $\overline{ac}$ and $\overline{ab}$, so
\[
(10b+c)+(10a+c)+(10a+b)=20a+11b+2c=293 .
\]
If $a\le8$, the left side is at most $160+99+18=277<293$, so $a=9$ and $11b+2c=113$. Then $11b=113-2c\ge113-18=95$ forces $b=9$, and $2c=14$, $c=7$. The number is $\boxed{997}$. Check: $97+97+99=293$.

(If the three digits were also required to be distinct, there would be no solution at all: the largest possible value of $20a+11b+2c$ with distinct digits is $180+88+14=282<293$.)
